\documentclass[11pt,a4paper]{article}
\usepackage[T1]{fontenc}
\usepackage[utf8]{inputenc}
\usepackage{lmodern}
\usepackage{microtype}
\usepackage[margin=29mm,headheight=15pt]{geometry}
\usepackage{amsmath,amssymb,amsthm,mathtools}
\usepackage{booktabs,array,longtable}
\usepackage{xcolor}
\usepackage{listings}
\usepackage{enumitem}
\usepackage{fancyhdr}
\usepackage{xurl}
\usepackage[hidelinks]{hyperref}
\usepackage[nameinlink,noabbrev]{cleveref}
\pagestyle{fancy}
\fancyhf{}
\fancyhead[L]{\small Four Motzkin congruences}
\fancyhead[R]{\small OEIS A001006}
\fancyfoot[C]{\thepage}
\setlength{\parskip}{4pt}
\setlength{\parindent}{1.2em}
\setcounter{tocdepth}{2}
\numberwithin{equation}{section}
\newtheorem{theorem}{Theorem}[section]
\newtheorem{lemma}[theorem]{Lemma}
\newtheorem{corollary}[theorem]{Corollary}
\newtheorem{proposition}[theorem]{Proposition}
\theoremstyle{definition}
\newtheorem{definition}[theorem]{Definition}
\newtheorem{example}[theorem]{Example}
\theoremstyle{remark}
\newtheorem{remark}[theorem]{Remark}
\DeclareMathOperator{\CT}{CT}
\DeclareMathOperator{\supp}{supp}
\newcommand{\Z}{\mathbb Z}
\newcommand{\F}{\mathbb F}
\newcommand{\eps}{\varepsilon}
\newcommand{\M}{M}
\newcommand{\oeis}[1]{\href{https://oeis.org/#1}{\texttt{#1}}}
\lstset{basicstyle=\ttfamily\small,columns=fullflexible,keepspaces=true,
  frame=single,framerule=0.3pt,breaklines=true,showstringspaces=false,
  xleftmargin=0.6em,xrightmargin=0.6em,aboveskip=0.8em,belowskip=0.8em}
\hypersetup{pdftitle={Four Motzkin Congruences: A Constant-Term Proof with Sharp Ranges},
 pdfauthor={Prepared for Vladimir Reshetnikov},
 pdfsubject={Proofs of the four Peter Bala conjectures recorded in OEIS A001006}}

\title{\textbf{Four Motzkin Congruences}\\[0.3em]
\large A Constant-Term Proof with Sharp Ranges}
\author{Prepared for Vladimir Reshetnikov}
\date{20 September 2026}

\begin{document}
\maketitle
\begin{abstract}
The OEIS entry for the Motzkin numbers records four congruences proposed by
Peter Bala on 10 February 2022. We give self-contained proofs of all four.
Two concern indices immediately preceding multiples of prime powers; two
assert repetition of an initial block modulo a prime or its square.
A single constant-term transfer theorem explains the block phenomena at
arbitrary prime-power precision. For the prime-square case, the arithmetic
input is obtained by factoring a quadratic over the sixth roots of unity
and proving a binomial coefficient congruence directly. We further show
that the two proposed initial-block ranges are optimal for every prime and
exponent in their respective statements, and compute the first nonzero
error explicitly. The boundary-index formulas are extended to all powers
of every odd prime. A digit algorithm evaluates a Motzkin number modulo an
odd prime without constructing the integer itself. Reproducible Python
programs accompany the article; 276,814 finite checks pass. The proofs are
mathematical arguments, not deductions from these computations. Related
results of Kohen and Pan--Sun are identified explicitly; no claim of
publication priority is inferred from an OEIS conjecture label.
\end{abstract}

\vspace{0.4em}
\noindent\textbf{Main outcome.} All four targeted formulas are true. In particular,
for every prime $p\geq5$ and $k\geq2$,
\[
 M_{p^k+n}\equiv M_n\pmod{p^2}
 \qquad(0\leq n\leq p^{k-1}-3),
\]
whereas at the very next index the difference is
\[
 M_{p^k+p^{k-1}-2}-M_{p^{k-1}-2}
 \equiv
 \begin{cases}-2p,&p\equiv1\pmod6,\\p,&p\equiv5\pmod6\end{cases}
 \pmod{p^2}.
\]
\newpage
{\small\setlength{\parskip}{0pt}
\tableofcontents
}
\newpage

\section{The sequence and the four targets}\label{sec:targets}

\subsection{Why the Motzkin numbers are a substantial target}
The Motzkin number $M_n$ counts paths from $(0,0)$ to $(n,0)$ whose steps
are $(1,1)$, $(1,0)$, and $(1,-1)$, and whose height never becomes negative.
The sequence starts
\[
 1,1,2,4,9,21,51,127,323,835,2188,5798,\ldots.
\]
It is \oeis{A001006}. Its many interpretations include noncrossing partial
matchings and restricted trees and permutations; the entry provides
references and numerous equivalent formulas~\cite{oeisM}. Thus the target
is a classical combinatorial sequence, not an isolated artificial example.

Our auxiliary sequences are the central and adjacent trinomial
coefficients. Put
\begin{equation}\label{eq:PQ}
 P(x)=x^{-1}+1+x,\qquad F(x)=1+x+x^2=xP(x),\qquad Q(x)=1-x^2,
\end{equation}
and define
\begin{equation}\label{eq:TU}
 T_n=\CT\,P(x)^n,\qquad U_n=[x^1]P(x)^n\qquad(n\geq0).
\end{equation}
Here $\CT$ extracts the coefficient of $x^0$ from a Laurent polynomial.
The $T_n$ form \oeis{A002426}, and the $U_n$ for $n\geq1$ form
\oeis{A005717}; we explicitly set $U_0=0$~\cite{oeisT,oeisU}.
Also write
\begin{equation}\label{eq:D}
 D_0=1,\qquad D_n=T_{n-1}+U_{n-1}\quad(n\geq1).
\end{equation}
This is \oeis{A005773}, the directed-animal counting sequence; the identity
in \eqref{eq:D} is recorded in that entry~\cite{oeisD}. The present proof
needs only its coefficient formulation.

\subsection{The exact statements to be proved}
The following are the four formulas in Bala's comment dated
10 February 2022 in \oeis{A001006}, with notation replaced by that just
introduced~\cite{oeisM}. The quantifier on $n$ is important, as are the
upper endpoints of the two ranges.

\begin{description}[style=nextline,leftmargin=1.5em,labelindent=0pt,labelwidth=0pt,font=\normalfont\bfseries]
\item[(B1) Boundary indices, first prime class.]
If $p\equiv1\pmod6$ is prime and $n,r\geq1$, then
\begin{equation}\tag{B1}\label{eq:B1}
 M_{np^r-2}\equiv-U_{n-1}\pmod p.
\end{equation}
\item[(B2) Boundary indices, second prime class.]
If $p\equiv5\pmod6$ is prime and $n\geq1$, then
\begin{equation}\tag{B2}\label{eq:B2}
 M_{np-2}\equiv-D_n\pmod p.
\end{equation}
\item[(B3) Repetition modulo a prime.]
If $p\geq3$ is prime and $k\geq1$, then
\begin{equation}\tag{B3}\label{eq:B3}
 M_{n+p^k}\equiv M_n\pmod p
 \qquad(0\leq n\leq p^k-3).
\end{equation}
\item[(B4) Repetition modulo the square of a prime.]
If $p\geq5$ is prime and $k\geq2$, then
\begin{equation}\tag{B4}\label{eq:B4}
 M_{n+p^k}\equiv M_n\pmod{p^2}
 \qquad(0\leq n\leq p^{k-1}-3).
\end{equation}
\end{description}

The proofs below resolve precisely these four assertions, not every
conjecture mentioned elsewhere in the Motzkin entry.

\subsection{Status, attribution, and what is established here}\label{sec:status}
The four formulas above were still grouped under a conjecture heading in
the OEIS entry consulted for this article. A database label is not a
reliable certificate that a statement is absent from the literature.
The entry itself includes a November 2024 comment by Rob Burns pointing
out that the $n=1$ boundary cases were already consequences of his
2017 work~\cite{oeisM}.
Kohen's 2024 work develops constant-term and base-$p$ methods for Motzkin
and central trinomial numbers~\cite{kohenUniform,kohenDensity}. In
particular, the prime-modulus formulas here overlap with consequences of
his digit decompositions and boundary identities. Pan and Sun establish
considerably stronger central-trinomial congruences than the prime-square
input needed here~\cite{panSun}; their paper first appeared as a preprint
in 2020 and its final version appeared in 2026.

The contribution of this note is a complete, directly checkable treatment
of the four OEIS-listed formulas together, including the extension to
higher prime powers, the optimality of both initial-block ranges, and the
explicit first errors. Every lemma needed for (B1)--(B4) is proved below.
Neither the proofs nor the claims of sharpness depend on a literature
search establishing novelty. No priority claim is made for the auxiliary
congruences or for consequences already implicit in the cited work.

\section{Constant terms and elementary identities}\label{sec:constant}

\subsection{A reflection argument}
\begin{lemma}\label{lem:reflection}
For all $n\geq0$,
\begin{equation}\label{eq:Mct}
 M_n=\CT\,Q(x)P(x)^n=[x^n](1-x^2)F(x)^n.
\end{equation}
\end{lemma}
\begin{proof}
The coefficient of $x^h$ in $P(x)^n$ counts unrestricted paths of length
$n$ ending at height $h$. Thus $T_n$ counts all paths ending at height
zero. A path of this kind that reaches height $-1$ has a first visit to
$-1$. Exchange its up-steps and down-steps through that visit and leave
all later steps unchanged. The transformed path ends at height $2$.
Conversely, any unrestricted path ending at height $2$ has a first visit
to height $1$; exchanging up and down through that visit recovers the bad
path. Horizontal steps are unchanged in both directions.

Consequently $M_n=[x^0]P^n-[x^2]P^n$. The polynomial $P^n$ is symmetric
under $x\mapsto x^{-1}$, so this difference is $\CT(1-x^2)P^n$.
Multiplication by $x^n$ gives the second expression.
\end{proof}

For later use observe the exact support bound
\begin{equation}\label{eq:support}
 \supp\bigl(Q(x)P(x)^n\bigr)\subseteq[-n,n+2]\cap\Z.
\end{equation}
It is the extra two units at the right edge, caused by $-x^2$, that
produce both the range cutoffs and their first corrections.

\subsection{The three neighboring coefficients}
Multiplication by $P(x)$ and symmetry give
\begin{equation}\label{eq:Uidentity}
 T_{n+1}=T_n+2U_n.
\end{equation}
Writing $V_n=[x^2]P(x)^n$, we also have
$T_{n+2}=3T_n+4U_n+2V_n$. Since $M_n=T_n-V_n$, elimination yields
\begin{equation}\label{eq:MTidentity}
 \boxed{\quad 2M_n=3T_n+2T_{n+1}-T_{n+2}.\quad}
\end{equation}
This is the identity that will turn the digit rule for $T_n$ into a fast
algorithm for $M_n$.

\subsection{Generating functions and exact recurrences}
Selecting the up- and down-steps of an unrestricted path gives
\begin{equation}\label{eq:Tsum}
 T_n=\sum_{j=0}^{\lfloor n/2\rfloor}\binom n{2j}\binom{2j}j.
\end{equation}
Consequently, as formal power series,
\begin{align}
 \mathcal T(z)=\sum_{n\geq0}T_nz^n
 &=\frac1{1-z}\sum_{j\geq0}\binom{2j}j
       \left(\frac{z^2}{(1-z)^2}\right)^j\notag\\
 &=\frac1{\sqrt{1-2z-3z^2}}.\label{eq:Tgf}
\end{align}
Differentiation gives
\begin{equation}\label{eq:Trec}
 (n+1)T_{n+1}=(2n+1)T_n+3nT_{n-1}\qquad(n\geq1).
\end{equation}

A nonempty Motzkin path either begins with a horizontal step, or has a
unique decomposition $U\pi D\rho$, where $D$ is the step of its first
return to height zero. Both $\pi$ and $\rho$ are Motzkin paths after a
vertical translation when needed. Therefore
\begin{equation}\label{eq:Mgf}
 \mathcal M(z)=1+z\mathcal M(z)+z^2\mathcal M(z)^2
 =\frac{1-z-\sqrt{1-2z-3z^2}}{2z^2}.
\end{equation}
Deleting horizontal steps from a path with $j$ up-steps leaves a Dyck
path of semilength $j$. Hence
\begin{equation}\label{eq:Msum}
 M_n=\sum_{j=0}^{\lfloor n/2\rfloor}\binom n{2j}
                  \frac1{j+1}\binom{2j}j.
\end{equation}
From \eqref{eq:Mgf}, differentiation and simplification give
\[
 z(1-2z-3z^2)\mathcal M'(z)
 +(2-3z-3z^2)\mathcal M(z)=2.
\]
Comparing coefficients yields the exact integer recurrence
\begin{equation}\label{eq:Mrec}
 (n+2)M_n=(2n+1)M_{n-1}+3(n-1)M_{n-2}\quad(n\geq2),
 \qquad M_0=M_1=1.
\end{equation}
These recurrences are useful for exact verification. In modular
arithmetic their denominators must not be inverted when divisible by the
modulus's prime; this point matters in the supplied implementation.

\section{Prime powers spread the exponents apart}\label{sec:transfer}

\subsection{Lifting a polynomial congruence}
All congruences between Laurent polynomials in this section are
coefficientwise. They take place in $\Z[x,x^{-1}]$.

\begin{lemma}\label{lem:lift}
Let $p$ be prime. If $A\equiv B\pmod{p^a}$ for $a\geq1$, then
$A^p\equiv B^p\pmod{p^{a+1}}$.
\end{lemma}
\begin{proof}
Write $A=B+p^aC$ and expand. The $j=1$ term is divisible by $p^{a+1}$.
For $1<j<p$, the factor $\binom pj$ contributes $p$, and the other
factor contributes $p^{aj}$, which is more than sufficient. The final
term is divisible by $p^{ap}$, and $ap\geq a+1$.
\end{proof}

\begin{lemma}\label{lem:spread}
For a Laurent polynomial $L$ with integer coefficients, a prime $p$, and
integers $1\leq s\leq k$, put $q=p^{k-s+1}$. Then
\begin{equation}\label{eq:spread}
 L(x)^{p^k}\equiv L(x^q)^{p^{s-1}}\pmod{p^s}.
\end{equation}
\end{lemma}
\begin{proof}
The Frobenius congruence in characteristic $p$ gives
$L(x)^q\equiv L(x^q)\pmod p$. Applying \cref{lem:lift} successively
$s-1$ times raises the exponent on both sides to $p^{s-1}$ and the
modulus to $p^s$.
\end{proof}

Notice the difference between precision and shift. To retain $s$ digits
of $p$-adic precision, we replace the large exponent $p^k$ by the smaller
exponent $p^{s-1}$, but the surviving exponents of $x$ are then spaced
by $q=p^{k-s+1}$. Only the coefficients near exponent zero can interact
with a sufficiently short factor $QP^n$.

\subsection{The full local-block theorem}
\begin{theorem}[Prime-power block transfer]\label{thm:block}
Let $p$ be an odd prime, $m\geq0$, and $1\leq s\leq k$. Set
\[
 q=p^{k-s+1},\qquad h=mp^{s-1}.
\]
Then, for $0\leq d<q$,
\begin{equation}\label{eq:block}
 M_{mp^k+d}\equiv
 \begin{cases}
 T_hM_d,&0\leq d\leq q-3,\\
 T_hM_{q-2}-U_h,&d=q-2,\\
 T_hM_{q-1}-(q-1)U_h,&d=q-1
 \end{cases}
 \pmod{p^s}.
\end{equation}
\end{theorem}
\begin{proof}
By \cref{lem:spread}, raising to the nonnegative integer power $m$ gives
\[
 P(x)^{mp^k}\equiv P(x^q)^h\pmod{p^s}.
\]
Thus \cref{lem:reflection} gives
\begin{equation}\label{eq:blockCT}
 M_{mp^k+d}\equiv\CT\bigl(Q(x)P(x)^dP(x^q)^h\bigr)\pmod{p^s}.
\end{equation}
Every exponent in the last factor is a multiple of $q$.

If $d\leq q-3$, the support of $QP^d$ lies strictly between $-q$ and
$q$, by \eqref{eq:support}. Only its constant term can meet an opposite
exponent from the last factor. The constant term of $P(x^q)^h$ is
$T_h$, and the first case follows.

If $d=q-2$, the support is contained in $[-q+2,q]$. Besides the constant
term, exactly one possible multiple of $q$ remains, namely $q$. Its
coefficient is $-1$: it comes from $-x^2$ times the top monomial
$x^{q-2}$ of $P^{q-2}$. The coefficient of $x^{-q}$ in $P(x^q)^h$ is
$U_h$ by symmetry. This gives the second case.

If $d=q-1$, the support is contained in $[-q+1,q+1]$. Since $q\geq3$,
again the only nonzero multiple of $q$ that can occur is $q$. Its
coefficient is
\[
 [x^q]QP^{q-1}=-[x^{q-2}]P^{q-1}=-(q-1).
\]
Indeed, an exponent one below the maximum is obtained by choosing the
constant monomial from exactly one of the $q-1$ factors, and $x$ from
all the others. Substitution in \eqref{eq:blockCT} gives the last case.
\end{proof}

\begin{remark}
The theorem is uniform in $m$ and contains all three residues near the
end of a block. It is not an assertion of global periodicity. In the
interior, the multiplying factor is a central trinomial coefficient; at
the boundary, an adjacent coefficient enters as well.
\end{remark}

\subsection{Proof and exact range of (B3)}
Taking $s=1$ and $m=1$ in \cref{thm:block} gives $q=p^k$, $h=1$, and
$T_1=U_1=1$. The interior case is exactly (B3). The next case gives
\begin{corollary}\label{cor:sharpP}
For every odd prime $p$ and $k\geq1$,
\begin{equation}\label{eq:sharpP}
 M_{2p^k-2}-M_{p^k-2}\equiv-1\pmod p.
\end{equation}
Consequently $p^k-3$ is the largest upper endpoint $L$ for which
$M_{p^k+n}\equiv M_n\pmod p$ holds for every $0\leq n\leq L$.
\end{corollary}
For example, when $p=5$ and $k=1$, the differences for $n=0,1,2$ vanish
modulo $5$, whereas $M_8-M_3=323-4=319\equiv-1\pmod5$.

\section{Boundary residues and proofs of (B1) and (B2)}\label{sec:boundary}

\subsection{An elementary period-three coefficient sequence}
For an odd prime $p$, define
\begin{equation}\label{eq:eps}
 \eps_p=
 \begin{cases}
 0,&p=3,\\
 1,&p\equiv1\pmod3,\\
 -1,&p\equiv2\pmod3.
 \end{cases}
\end{equation}
No quadratic reciprocity is needed for the following calculation.
We use only
\begin{equation}\label{eq:inverseF}
 \frac1{F(x)}=\frac{1-x}{1-x^3}=\sum_{j\geq0}c_jx^j,
 \qquad
 c_j=\begin{cases}1,&j\equiv0\pmod3,\\-1,&j\equiv1\pmod3,\\0,&j\equiv2\pmod3.
 \end{cases}
\end{equation}

\begin{lemma}\label{lem:baseboundary}
For every odd prime $p$, every $r\geq1$, and $q=p^r$,
\begin{align}
 T_{q-1}&\equiv\eps_p^r\pmod p,\label{eq:Tqminus1}\\
 M_{q-2}&\equiv\frac{\eps_p^r-1}{2}\pmod p,\label{eq:Mqminus2}\\
 M_{q-1}&\equiv\frac{1+3\eps_p^r}{2}\pmod p.\label{eq:Mqminus1}
\end{align}
The fractions mean multiplication by the inverse of $2$ modulo $p$.
\end{lemma}
\begin{proof}
Since $F(x)^q\equiv F(x^q)\pmod p$ and $F(0)=1$, division by $F$ or
$F^2$ is legitimate in $\F_p[[x]]$. At any degree below $q$, the
factor $F(x^q)=1+x^q+x^{2q}$ contributes only its constant term.
Therefore
\[
 T_{q-1}=[x^{q-1}]F^{q-1}\equiv[x^{q-1}]F^{-1}=c_{q-1}.
\]
If $p\ne3$, the residue of $q=p^r$ modulo $3$ is determined by the
parity of $r$ when $p\equiv2\pmod3$; the result is $\eps_p^r$.
If $p=3$, the coefficient is $c_{q-1}=0$. This proves
\eqref{eq:Tqminus1}.

For the second identity, \eqref{eq:Mct} gives
\[
 M_{q-2}\equiv[x^{q-2}]\frac{1-x^2}{F(x)^2}.
\]
The rational function here is a derivative:
\begin{equation}\label{eq:derivative}
 \frac{1-x^2}{F(x)^2}=\left(\frac{x}{F(x)}\right)'.
\end{equation}
Thus $M_{q-2}\equiv(q-1)c_{q-2}\equiv-c_{q-2}\pmod p$.
If $q\equiv1\pmod3$ this is zero; if $q\equiv2\pmod3$ it is $-1$.
For $p=3$ it is $1$, which equals $-1/2$ modulo $3$. These are precisely
\eqref{eq:Mqminus2}.

Finally,
\[
 M_{q-1}\equiv[x^{q-1}]\frac{1-x^2}{F(x)}=c_{q-1}-c_{q-3}.
\]
The three cases are respectively $2$, $-1$, and $-1$ when $q$ is
$1$, $2$, or $0$ modulo $3$. They give \eqref{eq:Mqminus1}, including
$p=3$ where $1/2\equiv-1\pmod3$.
\end{proof}

\subsection{A single formula covering every odd prime}
\begin{theorem}[Boundary-index formula]\label{thm:boundary}
For every odd prime $p$ and every $n,r\geq1$,
\begin{equation}\label{eq:unified}
 \boxed{\quad M_{np^r-2}\equiv
       \frac{\eps_p^r-1}{2}\,T_{n-1}-U_{n-1}\pmod p.\quad}
\end{equation}
\end{theorem}
\begin{proof}
Use \cref{thm:block} with precision $s=1$, exponent $k=r$, multiplier
$m=n-1$, and offset $d=p^r-2$. Its second case gives
\[
 M_{np^r-2}\equiv T_{n-1}M_{p^r-2}-U_{n-1}\pmod p.
\]
Now use \eqref{eq:Mqminus2}.
\end{proof}

If $p\equiv1\pmod6$, then $\eps_p=1$ and the first term vanishes.
This proves (B1), including $n=1$ because $U_0=0$.
If $p\equiv5\pmod6$ and $r=1$, then $\eps_p=-1$, so the right-hand
side is $-T_{n-1}-U_{n-1}=-D_n$. This proves (B2).

The same theorem provides a useful extension that is not present in the
wording of (B2):
\begin{equation}\label{eq:alternation}
 M_{np^r-2}\equiv
 \begin{cases}
 -D_n,&p\equiv5\pmod6,\ r\text{ odd},\\
 -U_{n-1},&p\equiv5\pmod6,\ r\text{ even}
 \end{cases}
 \pmod p.
\end{equation}
At $p=3$ it becomes
\begin{equation}\label{eq:p3boundary}
 M_{n3^r-2}\equiv T_{n-1}-U_{n-1}\pmod3.
\end{equation}
Thus simply replacing $p$ by $p^r$ on the left of (B2), without regard
to the parity of $r$, would give the wrong result.

\section{The prime-square arithmetic input}\label{sec:squareinput}

The transfer theorem at precision two leaves the scalar $T_{mp}$.
We now prove $T_{mp}\equiv T_m\pmod{p^2}$ for $p\geq5$. The proof
uses neither a classification of Motzkin automata nor a preexisting
supercongruence theorem. A related root-of-unity identity is central to
the stronger results of Pan and Sun~\cite{panSun}.

\subsection{A binomial coefficient congruence}
\begin{lemma}[Babbage congruence]\label{lem:babbage}
For every prime $p$, integer $m\geq0$, and $0\leq j\leq m$,
\begin{equation}\label{eq:babbage}
 \binom{mp}{jp}\equiv\binom mj\pmod{p^2}.
\end{equation}
Also $p\mid\binom{mp}{j}$ whenever $0\leq j\leq mp$ and $p\nmid j$.
\end{lemma}
\begin{proof}
There is an integer polynomial $E(x)$ supported on exponents
$1,\ldots,p-1$ such that
\[
 (1+x)^p=1+x^p+pE(x).
\]
For $m\geq1$, expansion to first order in $p$ gives
\[
 (1+x)^{mp}\equiv(1+x^p)^m+mpE(x)(1+x^p)^{m-1}\pmod{p^2}.
\]
The second term has no exponent divisible by $p$. Extracting $[x^{jp}]$
proves \eqref{eq:babbage}. The case $m=0$ is immediate.
For the final assertion use
$j\binom{mp}{j}=mp\binom{mp-1}{j-1}$ and invert $j$ modulo $p$.
\end{proof}

\subsection{Why the second power of the prime appears}
Let
\[
 \alpha=\frac{1+\sqrt{-3}}2,\qquad
 \beta=\frac{1-\sqrt{-3}}2.
\]
Then $\alpha+\beta=1$, $\alpha\beta=1$, and
$\alpha^6=\beta^6=1$. Moreover
\begin{equation}\label{eq:factor}
 F(x)=(x+\alpha)(x+\beta).
\end{equation}
Extracting the central coefficient of the $n$th power gives
\begin{equation}\label{eq:roots}
 T_n=\sum_{j=0}^n\binom nj^2\alpha^{n-j}\beta^j.
\end{equation}
The two binomial coefficients, arising from the two linear factors,
explain the extra $p$-divisibility.

To be precise about the ambient ring, work in
$R=\Z[\alpha]$, with $\beta=1-\alpha$. It is a free abelian group
with basis $1,\alpha$. In particular
\begin{equation}\label{eq:intersection}
 p^aR\cap\Z=p^a\Z.
\end{equation}
Indeed, the coefficient of $\alpha$ in an integer is zero, and its
coefficient of $1$ is that integer. Thus a congruence between integers
proved coefficientwise in this basis is an ordinary integer congruence.

\begin{theorem}\label{thm:Tp}
For all primes $p\geq5$ and integers $m\geq0$,
\begin{equation}\label{eq:Tp}
 T_{mp}\equiv T_m\pmod{p^2}.
\end{equation}
In particular $T_p\equiv1\pmod{p^2}$.
\end{theorem}
\begin{proof}
In \eqref{eq:roots} with $n=mp$, any term for which $p\nmid j$
vanishes modulo $p^2$, by \cref{lem:babbage}. For $j=ip$, the same
lemma allows replacement of $\binom{mp}{ip}^2$ by $\binom mi^2$.
Consequently
\[
 T_{mp}\equiv\sum_{i=0}^m\binom mi^2
                    (\alpha^p)^{m-i}(\beta^p)^i\pmod{p^2R}.
\]
Every prime $p\geq5$ is $1$ or $5$ modulo $6$. In the first case
$(\alpha^p,\beta^p)=(\alpha,\beta)$; in the second it is
$(\beta,\alpha)$. The sum is therefore $T_m$ in either case, using
$i\mapsto m-i$ in the second. Finally apply \eqref{eq:intersection}.
\end{proof}

\begin{remark}
The exclusion of $p=3$ has a visible algebraic cause: both
$\alpha^3$ and $\beta^3$ equal $-1$, rather than preserving or
exchanging the roots. Indeed $T_3=7\not\equiv1\pmod9$.
This is not a technical omission in the proof.
\end{remark}

\section{Proof of (B4), multiplier extension, and sharpness}\label{sec:B4}

\subsection{A stronger interior congruence}
\begin{theorem}\label{thm:multiplier}
For every prime $p\geq5$, $k\geq2$, and $m\geq0$,
\begin{equation}\label{eq:multiplier}
 \boxed{\quad M_{mp^k+n}\equiv T_mM_n\pmod{p^2}
 \quad(0\leq n\leq p^{k-1}-3).\quad}
\end{equation}
\end{theorem}
\begin{proof}
Set $s=2$ in \cref{thm:block}. Its interior term is
$T_{mp}M_n$, and \cref{thm:Tp} replaces $T_{mp}$ by $T_m$ modulo
$p^2$.
\end{proof}
Taking $m=1$ proves (B4), since $T_1=1$. This completes the proof of
all four OEIS-listed targets.

\subsection{The first adjacent coefficient modulo \texorpdfstring{$p^2$}{p squared}}
\begin{lemma}\label{lem:Up}
For $p\geq5$ prime,
\begin{equation}\label{eq:Up}
 U_p\equiv\frac p2(1+3\eps_p)\equiv
 \begin{cases}2p,&p\equiv1\pmod6,\\-p,&p\equiv5\pmod6\end{cases}
 \pmod{p^2}.
\end{equation}
More generally, for $m\geq1$,
\begin{equation}\label{eq:Ump}
 U_{mp}\equiv\frac{mp}{2}
            \bigl(T_m+3\eps_pT_{m-1}\bigr)\pmod{p^2}.
\end{equation}
\end{lemma}
\begin{proof}
Combining \eqref{eq:Trec} with \eqref{eq:Uidentity} gives the exact
identity
\begin{equation}\label{eq:Urec}
 2(N+1)U_N=N\bigl(T_N+3T_{N-1}\bigr)\qquad(N\geq1).
\end{equation}
For $N=mp$, the factor $N+1$ is invertible modulo $p^2$, and
$N/(N+1)\equiv N\pmod{p^2}$. It is enough to know the coefficients
inside the parentheses modulo $p$.

Theorem~\ref{thm:Tp} gives $T_{mp}\equiv T_m\pmod p$. Also,
by Frobenius,
\[
 T_{mp-1}=\CT\,P^{(m-1)p}P^{p-1}
 \equiv\CT\,P(x^p)^{m-1}P(x)^{p-1}
 =T_{m-1}T_{p-1}\pmod p.
\]
The last equality holds because the support of $P^{p-1}$ is strictly
between $-p$ and $p$. Equation~\eqref{eq:Tqminus1} gives
$T_{p-1}\equiv\eps_p$. Substituting in \eqref{eq:Urec} proves
\eqref{eq:Ump}. Its case $m=1$ is \eqref{eq:Up}.
\end{proof}

\subsection{The maximal initial interval}
\begin{theorem}[Exact first failure modulo $p^2$]\label{thm:sharpP2}
For every prime $p\geq5$ and integer $k\geq2$,
\begin{equation}\label{eq:sharpP2}
 \boxed{\quad
 M_{p^k+p^{k-1}-2}-M_{p^{k-1}-2}
 \equiv
 \begin{cases}
 -2p,&p\equiv1\pmod6,\\
 p,&p\equiv5\pmod6
 \end{cases}
 \pmod{p^2}.\quad}
\end{equation}
In particular $p^{k-1}-3$ is the largest upper endpoint $L$ for which
$M_{p^k+n}\equiv M_n\pmod{p^2}$ holds for every $0\leq n\leq L$.
\end{theorem}
\begin{proof}
Use the second case of \cref{thm:block} with $s=2$, $m=1$, and
$q=p^{k-1}$. It gives
\[
 M_{p^k+q-2}-M_{q-2}
 \equiv(T_p-1)M_{q-2}-U_p\equiv-U_p\pmod{p^2}.
\]
Now apply \cref{lem:Up}. Both resulting residues are nonzero modulo
$p^2$ because $p\geq5$. The preceding interior interval is valid by
\cref{thm:multiplier}, so its endpoint is maximal.
\end{proof}

The assertion is about the maximal interval \emph{starting at zero}.
Some later individual offsets may again satisfy the congruence; that does
not contradict the exact first failure.

\begin{example}
For $p=5$, $k=2$, (B4) applies to offsets $0,1,2$. At the next offset,
\[
 M_{28}=208023278209,\qquad M_3=4,
\]
so $M_{28}-M_3\equiv5\pmod{25}$. For $p=7$, $k=2$, the first
excluded offset is $5$, and the difference is $35$ modulo $49$, equivalently $-14$ modulo $49$.
The same defect $35$ occurs at offset $47$ for the shift $7^3$.
\end{example}

\section{Higher precision: the correct extension and its limitations}\label{sec:higher}

\subsection{A transfer formula, not unrestricted periodicity}
The interior of \cref{thm:block} already proves
\begin{equation}\label{eq:precision}
 M_{mp^k+n}\equiv T_{mp^{s-1}}M_n\pmod{p^s},
 \quad
 \begin{gathered}
 p\text{ odd prime},\quad 1\leq s\leq k,\quad m\geq0,\\
 0\leq n\leq p^{k-s+1}-3.
 \end{gathered}
\end{equation}
At $s=1$ the scalar is $T_m$, and at $s=2$ it can again be replaced
by $T_m$ for $p\geq5$. At higher precision it is generally incorrect
to replace it by $T_m$, even when $m=1$.

For instance, choose $p=5$, $s=k=3$, and $n=0$. The theorem says
\[
 M_{125}\equiv T_{25}\pmod{125}.
\]
Exact integer computation gives
\begin{equation}\label{eq:p3counter}
 M_{125}\equiv T_{25}\equiv51\not\equiv1=M_0\pmod{125}.
\end{equation}
Thus even the reduced range $0\leq n\leq p^{k-2}-3$ does not make
the unweighted analogue of (B4) true modulo $p^3$.
This is a counterexample to a \emph{possible strengthening}, not to any
of the four original statements.

The prime restriction in (B4) is also essential:
\begin{equation}\label{eq:p3failure}
 M_9=835\equiv7\not\equiv1=M_0\pmod9.
\end{equation}
The transfer theorem correctly predicts the factor $T_3=7$ here.

\subsection{Relationship to stronger known arithmetic}
Pan--Sun's theorem gives, for $p>3$ and $a\geq1$,
\[
 T_{mp^a}\equiv T_{mp^{a-1}}\pmod{p^{2a+2v_p(m)}}
 \qquad(m\geq1),
\]
as a consequence of the $p$-adic integrality of
$(T_{pN}-T_N)/(pN)^2$~\cite{panSun}. This is much stronger than
\cref{thm:Tp}. It can be inserted into \eqref{eq:precision} to reduce
its scalar further. For example, for $s\geq2$, put
$t=\lfloor(s-1)/2\rfloor$. All successive reductions from exponent
$s-1$ down to $t$ have modulus at least $p^s$, so
\begin{equation}\label{eq:panrefine}
 T_{mp^{s-1}}\equiv T_{mp^t}\pmod{p^s}.
\end{equation}
For $m=0$ the equality is immediate. This optional refinement depends
on the cited stronger theorem, unlike (B1)--(B4) and the block theorem,
whose proofs above are self-contained.

\section{A digit algorithm for enormous indices}\label{sec:algorithm}

\subsection{The Lucas rule for central trinomial coefficients}
\begin{proposition}\label{prop:Lucas}
For a prime $p$, $a\geq0$, and $0\leq b<p$,
\begin{equation}\label{eq:Lucas}
 T_{ap+b}\equiv T_aT_b\pmod p.
\end{equation}
Consequently, if $N=d_0+d_1p+\cdots+d_\ell p^\ell$ with
$0\leq d_i<p$, then
\begin{equation}\label{eq:digits}
 T_N\equiv\prod_{i=0}^{\ell}T_{d_i}\pmod p.
\end{equation}
\end{proposition}
\begin{proof}
Frobenius gives
$T_{ap+b}\equiv\CT\,P(x^p)^aP(x)^b\pmod p$.
All exponents of $P(x)^b$ lie in $[-b,b]\subset(-p,p)$, so only
constant terms can interact. This proves \eqref{eq:Lucas}, and
iteration gives \eqref{eq:digits}.
\end{proof}
This digit rule is also the foundation of the corresponding results in
Kohen's work~\cite{kohenUniform,kohenDensity}. Its proof is included so
that the algorithm's correctness does not rely on a black box.

\subsection{Evaluation and arithmetic complexity}
For an odd prime $p$, first compute $t_j=T_j\bmod p$ for $0\leq j<p$.
Recurrence~\eqref{eq:Trec} can be used here because its divisors are
strictly less than $p$. The inverses $1^{-1},\ldots,(p-1)^{-1}$ may be
computed in $O(p)$ field operations using
\[
 \operatorname{inv}(i)\equiv
 -\left\lfloor\frac pi\right\rfloor\operatorname{inv}(p\bmod i)\pmod p
 \qquad(2\leq i<p).
\]
This follows by dividing the identity
$p=\lfloor p/i\rfloor i+(p\bmod i)$ by $i(p\bmod i)$ in $\F_p$.

Use the base-$p$ digits of $N$, $N+1$, and $N+2$ to evaluate their
three central trinomial coefficients, and then apply
\begin{equation}\label{eq:Mfast}
 M_N\equiv2^{-1}\bigl(3T_N+2T_{N+1}-T_{N+2}\bigr)\pmod p.
\end{equation}
The table takes $O(p)$ storage and $O(p)$ field operations; each query
then takes $O(1+\log_p(N+2))$ field multiplications and integer digit
extractions. These are arithmetic-operation counts. For a growing input
integer $N$, the bit cost of repeated division to extract digits must
also be included; the statement is not a unit-cost bit-complexity bound.

A compact version of the evaluation stage is:
\begin{lstlisting}[language=Python]
def trinomial_from_digits(n, p, table):
    value = 1
    while n:
        n, digit = divmod(n, p)
        value = value * table[digit] % p
    return value


def motzkin_mod_prime(n, p, table):
    a = trinomial_from_digits(n, p, table)
    b = trinomial_from_digits(n + 1, p, table)
    c = trinomial_from_digits(n + 2, p, table)
    return (3*a + 2*b - c) * ((p + 1)//2) % p
\end{lstlisting}
The supplied module adds validation, construction of the table, exact
arithmetic routines, and a command-line interface.

\subsection{Prime-power block evaluation}
For an index of the form $mp^k+d$ covered by \cref{thm:block}, computation
modulo $p^s$ needs only $M_d$, $T_h$, and possibly $U_h$, where
$h=mp^{s-1}$. When $k$ is large relative to $s$, these can be dramatically
smaller than the target index. The implementation
\texttt{block\_mod\_prime\_power} uses precisely the three cases of
\eqref{eq:block}.

This is a specialized block algorithm, not a general fast algorithm for
arbitrary indices modulo arbitrary prime powers. Its local exact
computations grow with $d$ and $h$. Stating that restriction avoids
confusing a useful reduction with a universal complexity result.

\section{Reproducible checks and artifact contents}\label{sec:verification}

\subsection{Independent finite verification}
The accompanying verifier constructs $M_0,\ldots,M_{20000}$ and enough
$T_n$ values by the exact integer recurrences. Each division is checked
for a zero remainder before its result is used. It does not generate
these values using the congruences under test.

The first 151 values are independently checked against the binomial
sums \eqref{eq:Tsum} and \eqref{eq:Msum}; the first 81 Motzkin numbers
are also checked by height-by-height path dynamic programming. The
congruence checks cover odd primes up to $97$, with parameters restricted
so that the exact oracle contains the required indices. Multiplier
bounds and the complete iteration limits are explicit in the script.
The higher-precision block checks use $p\in\{3,5,7,11,13\}$,
$s\leq4$, and the same index cap.

\begin{center}
\begin{tabular}{@{}lr@{}}
\toprule
Check family & Number of checks\\
\midrule
Bala (B1) & 251\\
Bala (B2) & 144\\
Bala (B3) & 98,471\\
Bala (B4) & 10,476\\
Full prime-power block theorem & 96,478\\
Prime-modulus digit algorithm & 28,824\\
Three-trinomial identity & 20,001\\
Prime-square multiplier extension & 12,735\\
All remaining checks & 9,434\\
\midrule
Total & 276,814\\
\bottomrule
\end{tabular}
\end{center}

The exact first failure modulo $p^2$ was checked for 36 admissible
prime--exponent pairs. A few rows of the full CSV are:
\begin{center}
\begin{tabular}{@{}rrrrr@{}}
\toprule
$p$ & $k$ & First excluded offset & Observed error & Modulus\\
\midrule
5 & 2 & 3 & 5 & 25\\
5 & 3 & 23 & 5 & 25\\
7 & 2 & 5 & 35 & 49\\
7 & 3 & 47 & 35 & 49\\
11 & 2 & 9 & 11 & 121\\
13 & 2 & 11 & 143 & 169\\
\bottomrule
\end{tabular}
\end{center}
These observations are consequences of \cref{thm:sharpP2}, not the
basis of its proof. The programs are not a formal proof-assistant
verification of the mathematical arguments.

\subsection{Example commands}
The code uses only the Python standard library and is written for
Python 3.9 or later. The included run was executed with Python 3.13.5.
From the archive's top-level directory:
\begin{lstlisting}
python code/verify.py
python code/motzkin.py exact 28
python code/motzkin.py mod 12345678901234567890 97
python code/motzkin.py block 1 3 0 5 3
\end{lstlisting}
The last command returns $51$, illustrating \eqref{eq:p3counter}.
The command
\begin{lstlisting}
python code/verify.py --max-index 20000 --prime-limit 97
\end{lstlisting}
recreates the checked JSON report and CSV data. Increasing the exact
index cap consumes more memory because the verifier deliberately stores
large exact integers; this is verification infrastructure rather than the
space-efficient query algorithm.

For the much larger index $N=10^{100}+123456789$, the digit evaluator
returns the following residues:
\begin{center}
\begin{tabular}{@{}rrrrr@{}}
\toprule
$p$ & 5 & 7 & 11 & 97\\
\midrule
$M_N\bmod p$ & 1 & 0 & 1 & 28\\
\bottomrule
\end{tabular}
\end{center}
These last outputs illustrate the proved algorithm. They were not
independently compared with a construction of the enormous exact $M_N$,
and are not included as independent oracle checks.

\subsection{Files supplied}
The archive contains the \LaTeX{} source and its compiled PDF,
\nolinkurl{code/motzkin.py}, \nolinkurl{code/verify.py},
\nolinkurl{data/verification.json}, \nolinkurl{data/sharp_boundaries.csv},
and \nolinkurl{data/initial_values.csv}. The last file contains values
through index $200$, with sequence IDs and the $U_0=0$ convention made
explicit. A README gives build and execution instructions. Notes record
the source locations and a proposed OEIS update; nothing has been
submitted to or changed in OEIS.

\section{Conclusion}\label{sec:conclusion}
All four of Bala's formulas follow from a small set of exact mechanisms.
The factor $1-x^2$ imposes the nonnegativity condition on paths. Frobenius
and its elementary lifting spread the exponents far apart. In the
interior of a block, only the central coefficient survives. At its last
two offsets, a single adjacent coefficient survives as well. Finally,
factorization over sixth roots of unity turns prime divisibility of a
binomial coefficient into prime-square divisibility of its square.

The same reasoning settles the two natural questions that the original
congruences leave implicit. The prime-$5\bmod6$ boundary formula
alternates with the parity of the exponent, and both stated repetition
ranges end exactly where the first nonzero support interaction occurs.
The errors are explicit and never vanish in the relevant moduli:
$-1$ modulo $p$, and either $-2p$ or $p$ modulo $p^2$.
The higher-precision theorem retains the necessary central-trinomial
factor rather than asserting an invalid stronger periodicity.

\clearpage
\appendix
\section{An alternative rational proof of \texorpdfstring{$T_p\equiv1\pmod{p^2}$}{the scalar congruence}}
\label{app:rational}
The root-of-unity proof gives the useful multiplier statement
$T_{mp}\equiv T_m\pmod{p^2}$. For readers preferring a proof entirely
in rational arithmetic, the special case needed for (B4) has the
following independent derivation.

Let $L_0=2$, $L_1=1$, and $L_n=L_{n-1}-L_{n-2}$. The sequence is
periodic with period six:
\[
 2,1,-1,-2,-1,1,2,1,\ldots.
\]
In particular $L_p=1$ for every prime $p\geq5$. Its generating
identity is
\[
 -\log(1-z+z^2)=\sum_{n\geq1}\frac{L_n}{n}z^n.
\]
This can be checked by differentiation and the initial values. Expanding
$-\log(1-(z-z^2))$ and extracting coefficients gives, for $n\geq1$,
\begin{equation}\label{eq:Lucaspoly}
 L_n=\sum_{j=0}^{\lfloor n/2\rfloor}
 (-1)^j\frac{n}{n-j}\binom{n-j}{j}.
\end{equation}
Take $n=p$. For $1\leq j\leq(p-1)/2$,
\[
 \binom p{2j}\binom{2j}j
 =\frac p{2j}\binom{p-1}{2j-1}\binom{2j}j
 \equiv-\frac p{2j}\binom{2j}j\pmod{p^2}.
\]
On the other hand, the $j$th summand of \eqref{eq:Lucaspoly} is
\begin{align*}
 (-1)^j\frac p{p-j}\binom{p-j}j
 &=(-1)^j\frac pj\binom{p-j-1}{j-1}\\
 &\equiv-\frac pj\binom{2j-1}{j-1}
 =-\frac p{2j}\binom{2j}j\pmod{p^2}.
\end{align*}
All denominators in these congruences are prime to $p$. The $j=0$
summand in both expressions is $1$. Comparing with \eqref{eq:Tsum}
therefore proves $T_p\equiv L_p=1\pmod{p^2}$. This also explains
why the period-six distinction appears without requiring arithmetic in
an extension ring.

\clearpage
\phantomsection
\addcontentsline{toc}{section}{References}
\begin{thebibliography}{99}
\bibitem{oeisM}
The OEIS Foundation Inc., \emph{A001006: Motzkin numbers},
\url{https://oeis.org/A001006}. Consulted 20 September 2026.
The four target conjectures are the comment by Peter Bala dated
10 February 2022.

\bibitem{oeisT}
The OEIS Foundation Inc., \emph{A002426: Central trinomial coefficients},
\url{https://oeis.org/A002426}. Consulted 20 September 2026.

\bibitem{oeisU}
The OEIS Foundation Inc., \emph{A005717: Next-to-central trinomial
coefficients}, \url{https://oeis.org/A005717}.
Consulted 20 September 2026. The present article extends the offset by
setting $U_0=0$.

\bibitem{oeisD}
The OEIS Foundation Inc., \emph{A005773: Number of directed animals of
size $n$}, \url{https://oeis.org/A005773}.
Consulted 20 September 2026. See the coefficient identity in David
Callan's comment dated 7 February 2004.

\bibitem{kohenUniform}
N.~Kohen, \emph{Uniform Recurrence in the Motzkin Numbers and Related
Sequences mod $p$}, arXiv:2403.00149v1 (2024).
\url{https://arxiv.org/abs/2403.00149}.
See the central-trinomial digit rule and the use of
$2M_n=3T_n+2T_{n+1}-T_{n+2}$.

\bibitem{kohenDensity}
N.~Kohen, \emph{Density and Symmetry in the Generalized Motzkin Numbers
mod $p$}, arXiv:2411.03681v1 (2024).
\url{https://arxiv.org/abs/2411.03681}.
See Proposition 6 for $M_{p-2}$, Proposition 7 for digit
multiplicativity, and Proposition 9 for boundary digit patterns.

\bibitem{panSun}
H.~Pan and Z.-W.~Sun, \emph{Supercongruences for central trinomial
coefficients}, Frontiers in Combinatorics and Number Theory
\textbf{2} (2026), 55--62.
\url{https://arxiv.org/abs/2012.05121v3}.
First preprint posted 9 December 2020; final arXiv revision
13 March 2026. See Theorem 1.1 and Lemma 2.1.
\end{thebibliography}
\end{document}
